\chapter*{开始之前}
\addcontentsline{toc}{chapter}{开始之前}
\markboth{开始之前}{}

这本练习册带你把 Dublin Bikes 的“车站”功能用 Spring Boot 重写一遍，边做题边学。题目格式模仿 OA：每一关分几个小题，每个小题有说明、示例和一组测试，测试全绿就算过关。

\section*{做完以后你会得到什么}

\begin{center}
\begin{tikzpicture}[
  box/.style={draw=accent, fill=taskbg, rounded corners=2pt, align=center, font=\footnotesize,
              minimum height=1.3cm, text width=2.85cm, inner sep=3pt},
  ext/.style={box, draw=black!40, fill=examplebg},
  arr/.style={-{Stealth[length=2.2mm]}, thick, accent!80!black},
  node distance=0.45cm]
  \node[ext] (api) {JCDecaux API};
  \node[box, right=of api] (client) {\texttt{JcDecauxClient}\\[1pt]\scriptsize Part 3};
  \node[box, right=of client] (job) {\texttt{StationScrapeJob}\\[1pt]\scriptsize 每 5 分钟 · Part 5};
  \node[box, right=of job] (ingest) {\texttt{StationIngest-}\\\texttt{Service}\\[1pt]\scriptsize Part 5};
  \node[box, below=1.1cm of ingest] (db) {数据库\\[1pt]\scriptsize Repository · Part 4};
  \node[box, left=of db] (query) {\texttt{StationQuery-}\\\texttt{Service}\\[1pt]\scriptsize Part 6};
  \node[box, left=of query] (ctrl) {\texttt{StationController}\\[1pt]\scriptsize Part 6};
  \node[ext, left=of ctrl] (map) {浏览器地图\\[1pt]\scriptsize 已提供};
  \draw[arr] (client) -- (api);
  \draw[arr] (job) -- (client);
  \draw[arr] (job) -- (ingest);
  \draw[arr] (ingest) -- (db);
  \draw[arr] (query) -- (db);
  \draw[arr] (ctrl) -- (query);
  \draw[arr] (map) -- (ctrl);
\end{tikzpicture}
\end{center}

一个能真正运行的后端：定时抓取都柏林 113 个站点的实时数据，存进数据库，再通过 REST API 提供给地图页面。

\section*{关卡}

\begin{center}
\begin{tabularx}{\textwidth}{l X l l}
\toprule
关卡 & 内容 & 测试 & 预计用时 \\
\midrule
Part 1 & 现代 Java：record、Stream、Optional（不涉及 Spring） & \code{Part1Modern*} & 1 小时 \\
Part 2 & Spring 核心：容器、Bean、依赖注入、配置绑定 & \code{Part2*} & 1 小时 \\
Part 3 & 调用外部 API：RestClient、JSON 映射、错误处理 & \code{Part3*} & 1 小时 \\
Part 4 & 数据库：JPA 实体、复合主键、Repository 查询 & \code{Part4*} & 1.5--2 小时 \\
Part 5 & 事务和定时任务：\code{@Transactional}、\code{@Scheduled}、Mockito & \code{Part5*} & 1 小时 \\
Part 6 & REST API：Controller、参数校验、统一错误处理 & \code{Part6*} & 1.5 小时 \\
\midrule
Part 7 & 调用 Python 推理服务：POST、\code{@JsonNaming}、超时 & \code{Part7*} & 1 小时 \\
Part 8 & 重试、熔断、降级：Resilience4j & \code{Part8*} & 1.5 小时 \\
Part 9 & 预测接口：Clock、降级响应、两个方向的 DTO & \code{Part9*} & 1 小时 \\
Part 10 & API 文档与健康检查：springdoc、Actuator & \code{Part10*} & 1 小时 \\
\midrule
Part 11 & OpenWeather 客户端与缓存：\code{@Cacheable} & \code{Part11*} & 1 小时 \\
Part 12 & 天气入库：Flyway V2、\code{Limit}、两种写入策略 & \code{Part12*} & 1 小时 \\
Part 13 & 天气接口、多个 advice、时间匹配 & \code{Part13*} & 1 小时 \\
\midrule
Part 14 & 声明式 HTTP 客户端：\code{@HttpExchange} 调 Google Maps & \code{Part14*} & 1 小时 \\
Part 15 & 路线接口：\code{@RequestBody}、\code{@Valid}、按原因映射状态码 & \code{Part15*} & 1 小时 \\
\midrule
Part 16 & 调用 LLM：请求头认证、429 限流 & \code{Part16*} & 1 小时 \\
Part 17 & 聊天接口：提示词、不可信输入、\code{ResponseEntity} & \code{Part17*} & 1 小时 \\
\bottomrule
\end{tabularx}
\end{center}

\textbf{Part 7--17（ML 集成、工程化、天气、路线、聊天机器人）是第二批}，前提是 Part 1--6 已经全部做完。

\textbf{按顺序做。}后面的关卡会用到前面写的代码，但每一关的测试是独立的：做完 Part 2 时，Part 1、2 的测试会全绿，Part 3 以后的测试还是红的，这是正常的。

\section*{准备}

\begin{enumerate}
  \item 用 IntelliJ IDEA 打开仓库的根目录（有 \code{pom.xml} 的那一层），等它自动导入 Maven 项目。
  \item 确认 Project SDK 是 Java 25（File → Project Structure → Project）。
  \item 跑一次全部测试，确认环境没问题。应该能编译，每做完一关，那一关的测试会变绿（Part 1--6 做完时 58 个全绿；加上 Part 7--17 一共 152 个）：
\end{enumerate}

\begin{shell}
./mvnw test
\end{shell}

不需要安装 Maven，\code{mvnw} 会自动下载。不需要数据库，测试用的是内存数据库 H2。

\section*{做题流程}

\begin{enumerate}
  \item 读这一关的“知识点”，再读题目。
  \item 在 \code{src/main/java} 里搜 \code{TODO 1.1} 这样的标记，找到要改的地方。
  \item 跑这一关的测试：\code{./mvnw test -Dtest='Part2*'}（Part 1 要写成 \code{'Part1Modern*'}，因为 \code{Part1*} 也会匹配到 Part 10--13）。在 IntelliJ 里也可以直接点测试类旁边的绿色三角形，结果更直观。
  \item 测试全绿后提交一次 commit，再进入下一关。
\end{enumerate}

\section*{规则}

\begin{itemize}
  \item \textbf{不要改测试文件}（\code{src/test} 下面的所有文件）。测试就是题目要求。
  \item \textbf{不要改已经写好的类}：\code{NearbyStation}、\code{IngestResult}、\code{JcDecauxException}、\code{StationNotFoundException}、\code{StationResponse}、\code{NearbyStationResponse}，以及 Flyway 的 SQL 文件和 \code{application.yml}。
  \item 方法签名（方法名、参数、返回类型）不要改，否则测试编译不过。
  \item 卡住超过 20 分钟：先看附录里的提示，再问 AI 要思路。实在做不出来再看答案。
\end{itemize}

\section*{答案}

\code{exercises/solutions/} 里是每个文件的参考答案，路径和 \code{src/main/java/ie/dublinbikes/} 一一对应。答案已经验证过：全部 58 个测试通过，应用能用真实的 JCDecaux 数据启动。\textbf{建议做完一关再看答案}，对比一下自己的写法哪里不同。

\section*{全部做完以后}

\begin{shell}
./mvnw test                                  # 全部测试
SCRAPER_ENABLED=true ./mvnw spring-boot:run
\end{shell}

几秒后打开 \url{http://localhost:8080/map.html} 就能看到地图上的 113 个站点。运行前先在仓库根目录建一个 \code{.env} 文件，写一行 \code{JCDECAUX_API_KEY=你的key}（可以从原项目的 \code{var.env} 里复制）。\code{.env} 已经在 \code{.gitignore} 里，不会被提交。

想连真实的 MySQL 而不是内存数据库（这条路径还没有实际测试过）：

\begin{shell}
docker compose up -d
SCRAPER_ENABLED=true ./mvnw spring-boot:run -Dspring-boot.run.profiles=mysql
\end{shell}

\section*{用到的版本}

Java 25、Spring Boot 4.1.2（Spring Framework 7、Hibernate 7、Jackson 3）。网上很多教程是 Spring Boot 2 或 3 的，大部分写法一样，个别包名不同，以这里的代码为准。
